%%%%%
%
\section{Correction de valeurs hors-bornes avec flag}

\itemstyle{Op\'erateur~:} {\sf Correc\_hors\_bornes\_flag}
\index{Correc\_hors\_bornes\_flag}

%%%%%OBJET%%%%
\itemstyle{Objet~:}
Cette fonction d�tecte et remplace les valeurs hors d'un intervalle
sp�cifi�. Si la dur�e des valeurs fausses ou manquantes cons�cutives est courte
(inf�rieure � un seuil fix�), elles sont remplac�es par interpolation
lin�aire. Dans le cas contraire, la valeur en sortie est manquante.

Une variable flag indique si la valeur d'origine est valide,
manquante ou hors bornes.

%%%%%ARGUMENTS%%%%
\bigskip
%%%%%ENTREES%%%%
\begin{tabarg}
\multicolumn{3}{l}{\titretab {Entr\'ees}}\\
\hline
$X$ & Vecteur, \texttt{float} & param�tre\\
$borne_{inf}$ & Coeff., \texttt{float} & valeur min. admise\\
$borne_{sup}$ & Coeff., \texttt{float} & valeur max. admise\\
$seuil_{duree}$ & Coeff., \texttt{int} & Nombre maxi de valeurs cons�cutives\\
$MisVal$ & Coeff., \texttt{float} & valeur sp�ciale pour donn�e absente\\
\end{tabarg}

\medskip
%%%%%SORTIES%%%%
\begin{tabarg}
\multicolumn{3}{l}{\titretab {Sorties}}\\
\hline
$flag$ & Vecteur, \texttt{float} & flag en fonction de la valeur du parametre $X$ [-10, -1, 0, 1]\\
$X_{cor}$ & Vecteur, \texttt{float} & param�tre corrig� (m�me unit� que $X$)\\
\end{tabarg}


%%%%%FORMULE%%%%
\itemstyle{Formule}

%%Correc_hors_bornes_flag
1- recherche de p�riodes de valeurs hors-bornes cons�cutives.\newline
2- Pour chaque p�riode :\newline
\indent Si elle est de dur�e $< seuil_{duree}$, et encadr�e par 
des valeurs correctes,\newline
\hspace*{2cm} remplacement par une interpolation lin�aire\newline
sinon\newline
\hspace*{2cm} remplacement par des �valeurs manquantes�.\newline
\newline
En sortie : le flag permet de conna�tre l'origine des valeurs :
\begin{description}
\item[ ~0] si la valeur initiale $X$ est valide (non manquante et comprise entre les 2 bornes fix�es)
\item[ -1] si la valeur initiale $X$ est inf�rieure � la borne minimale
\item[ +1] si la valeur initiale $X$ est sup�rieure � la borne maximale
\item[-10] si la valeur initiale est manquante
\end{description}



%%%%%COMMENTAIRE%%%%
\itemstyle{Commentaire~:}
\begin{itemize}
\item Cette fonction est identique � Correc\_hors\_bornes avec l'ajout
d'un flag en sortie.
\item L'interpolation lin�aire ne convient pas aux variables
qui peuvent �tre naturellement discontinues, tels les caps et
autres angles.
\item Le flag peut s'utiliser de la fa�on suivante : si le flag
est nul, la valeur n'a pas �t� modifi�e (non interpol�e et non manquante).
\end{itemize}

%%%%%AUTEUR%%%%
\itemstyle{Auteur~:}
CNRM/GMEI/TRAMM

%%%%%REFERENCE%%%%
\itemstyle{R\'ef\'erences~:}

